%%%%%%%%%%%%%%%%%%%%%%%%%%%%%%%%%%%%%%%%%%%%%%%%%%%%%%%%%%%%%%%%%%%
% Vince's Tikz macros                                             %
% Version of 19 June 2016                                         %
%%%%%%%%%%%%%%%%%%%%%%%%%%%%%%%%%%%%%%%%%%%%%%%%%%%%%%%%%%%%%%%%%%%

\usepackage{pgfmath}

%
%
%
%

% Points of absolute size (to be used when drawing matchings and permutations).
% Call \absdot[label]{(2,3)} for label, otherwise omit optional argument.

\newcommand\mybullet{\raisebox{-5pt}{\normalsize \ensuremath{\bullet}}}
\newcommand\mycirc{\raisebox{-5pt}{\normalsize \ensuremath{\circ}}}
\newcommand\myscaledbullet[1]{\raisebox{-5pt}{\scalebox{#1}{\normalsize \ensuremath{\bullet}}}}

\makeatletter
\def\absdot{\@ifnextchar[{\@absdotlabel}{\@absdotnolabel}}
    \def\@absdotlabel[#1]#2{%
        \node at #2 {\normalsize \mybullet};
        \node at #2 [below=2pt] {\ensuremath{#1}};
    }
    \def\@absdotnolabel#1{%
        \node at #1 {\normalsize \mybullet};
    }
\def\absdothollow{\@ifnextchar[{\@absdothollowlabel}{\@absdothollownolabel}}
    % Make a dot of fixed absolute size.
    % Note that \absdothollow first overwrites what it is going on top of.
    \def\@absdothollowlabel[#1]#2{%
        \node at #2 {\normalsize \textcolor{white}{\mybullet}};
        \node at #2 {\normalsize \mycirc};
        \node at #2 [below=2pt] {\ensuremath{#1}};
    }
    \def\@absdothollownolabel#1{%
        \node at #1 {\normalsize \textcolor{white}{\mybullet}};
        \node at #1 {\normalsize \mycirc};
    }
\makeatother

%
%
%
%

% Plotting permutations. Recommended scale: 0.25 (and possibly up to 0.5).

\newcommand{\plotperm}[1]{
    \foreach \j [count=\i] in {#1} {
        \absdot{(\i,\j)};
    };
}

\newcommand{\plotpermprecise}[2]{
    \foreach \j [count=\i] in {#1} {
        \draw[black,fill=black] (\i,\j) circle (#2);
    };
}

\newcommand{\plotpermscaled}[2]{
    \foreach \j [count=\i] in {#1} {
        \node at (\i,\j) {\scalebox{#2}{\normalsize \ensuremath{\bullet}}};
    };
}

\newcommand{\plotpermhollow}[1]{
    \foreach \j [count=\i] in {#1} {
        \absdothollow{(\i,\j)};
    };
}

\newcommand{\plotpartialperm}[1]{
    \foreach \i/\j in {#1} {
        \absdot{(\i,\j)};
    };
}

\newcommand\myencircle{\raisebox{-8.75pt}{\LARGE \ensuremath{\circ}}}
\newcommand{\plotpartialpermencircle}[1]{
    \foreach \i/\j in {#1} {
        \node at (\i,\j) {\myencircle};
    };
}

\newcommand{\plotpartialpermhollow}[1]{
    \foreach \i/\j in {#1} {
        \absdothollow{(\i,\j)};
    };
}

\newcommand{\plotpermbox}[4]{
    \draw [darkgray, thick, line cap=round]
        ({#1-0.5}, {#2-0.5}) rectangle ({#3+0.5}, {#4+0.5});
}

\newcommand{\plotpermborder}[1]{
    \plotperm{#1};
    \pgfmathsetmacro\n{dim(#1)};
    % Now \n stores the number of entries of the permutation. Draw the border.
    \plotpermbox{1}{1}{\n}{\n};
}

\newcommand{\plotpermgraphedges}[1]{
    \foreach \j [count=\i] in {#1} {
        \foreach \b [count=\a] in {#1} {
            % Draw edge from (a,b) to (i,j) if they form an inversion.
            \ifthenelse{\a<\i \AND \b>\j}{\draw (\a,\b)--(\i,\j);}{}
        };
    };
}

\newcommand{\plotpermgraph}[1]{
    \foreach \j [count=\i] in {#1} {
        \foreach \b [count=\a] in {#1} {
            % Draw edge from (a,b) to (i,j) if they form an inversion.
            \ifthenelse{\a<\i \AND \b>\j}{\draw (\a,\b)--(\i,\j);}{}
        };
    };
    \plotperm{#1};
}

\newcommand{\plotpermgraphborder}[1]{
    \foreach \j [count=\i] in {#1} {
        \foreach \b [count=\a] in {#1} {
            % Draw edge from (a,b) to (i,j) if they form an inversion.
            \ifthenelse{\a<\i \AND \b>\j}{\draw (\a,\b)--(\i,\j);}{}
        };
    };
    \plotpermborder{#1};
}

\newcommand{\plotpermgrid}[1]{
    \pgfmathsetmacro\n{dim(#1)};
    % Now \n stores the number of entries of the permutation. Draw the border.
    \foreach \i in {0,1,2,...,\n}{
        \draw[color=darkgray] ({\i+0.5}, 0.5)--({\i+0.5}, {\n+0.5});
        \draw[color=darkgray] (0.5, {\i+0.5})--({\n+0.5}, {\i+0.5});
    }
    \plotperm{#1};
}

\newcommand{\plotpermdyckgrid}[1]{
    \pgfmathsetmacro\n{dim(#1)};
    % Now \n stores the number of entries of the permutation. Draw the border.
    \foreach \i in {0,1,2,...,\n}{
        \draw[color=lightgray] ({\i+0.5}, 0.5)--({\i+0.5}, {\n+0.5});
        \draw[color=lightgray] (0.5, {\i+0.5})--({\n+0.5}, {\i+0.5});
    }
    \draw[ultra thick, line cap=round, color=darkgray] (0.5, 0.5)--({\n+0.5}, {\n+0.5});
    \plotperm{#1};
}

\newcommand{\plotpermdyckpath}[1]{
    % Pass a list of coordinates. Everything is shifted, so the lower-left of the grid is the origin.
    \draw[ultra thick, line cap=round] (0.5,0.5)
    \foreach \step in {#1} {
        \ifnum\step=1
            -- ++(0,1)
        \else
            -- ++(1,0)
        \fi
    };
}

%
%
%
%

% Drawing Dyck paths. Pass the command a list of 1s and -1s.
% Note that this path will be a factor of sqrt(2) larger than \plotpermdyckpath
% Adapted from http://tex.stackexchange.com/questions/63540/how-to-draw-a-catalan-number-diagram-on-tikz

\newcommand{\plotdyckpath}[1]{
    \draw[ultra thick, line cap=round] (0.5,0)
    \foreach \step in {#1} {
        \ifnum\step=1
            -- ++(1,1)
        \else
            -- ++(1,-1)
        \fi
    };
}

%
%
%
%

% Drawing (large) matchings. Recommended scale is around 0.5.

\newcommand{\arcskinnyplain}[2]{
    \draw[thick] (#1,0) arc (180:0:{(#2-#1)/2});
}

\newcommand{\arc}[2]{
    \draw[thick] (#1,0) arc (180:0:{(#2-#1)/2});
    \absdot{(#1,0)};
    \absdot{(#2,0)};
}

\newcommand{\arcgray}[2]{
    \draw[thick, lightgray] (#1,0) arc (180:0:{(#2-#1)/2});
    \absdot{(#1,0)};
    \absdot{(#2,0)};
}

\newcommand{\arcdotted}[2]{
    \draw[thick, dotted, line cap=round] (#1,0) arc (180:0:{(#2-#1)/2});
    \absdot{(#1,0)};
    \absdot{(#2,0)};
}

\newcommand{\arclabel}[4]{
    \draw[thick] (#1,0) arc (180:0:{(#2-#1)/2});
    % The \mathstrut below gets the vertical alignment correct if you
    %   are using labels like a, b, c, d, e, ...
    \absdot[\ensuremath{\mathstrut #3}]{(#1,0)};
    \absdot[\ensuremath{\mathstrut #4}]{(#2,0)};
}

\newcommand{\arclabelgray}[4]{
    \draw[thick, darkgray] (#1,0) arc (180:0:{(#2-#1)/2});
    \absdot[\ensuremath{\mathstrut #3}]{(#1,0)};
    \absdot[\ensuremath{\mathstrut #4}]{(#2,0)};
}

\newcommand{\arclabeldotted}[4]{
    \draw[thick, dotted, line cap=round] (#1,0) arc (180:0:{(#2-#1)/2});
    \absdot[\ensuremath{\mathstrut #3}]{(#1,0)};
    \absdot[\ensuremath{\mathstrut #4}]{(#2,0)};
}

\newcommand{\arclabelgroup}[2]{
    \draw ({#1-0.25}, -1) rectangle ({#2+0.25}, -0.25);
}

\newcommand{\arclabelgroupgray}[2]{
    \draw[lightgray, fill=lightgray] ({#1-0.25}, -1) rectangle ({#2+0.25}, -0.3);
}

%
%
%
%

\newcommand{\matching}[1]{
    \foreach \i/\j in {#1} {
        \arc{\i}{\j};
    };
}

\newcommand{\matchingdotted}[1]{
    \foreach \i/\j in {#1} {
        \arcdotted{\i}{\j};
    };
}

\newcommand{\matchinglabels}[1]{
    \foreach \i/\j in {#1} {
        \arclabel{\i}{\j}{\i}{\j};
    };
}

\newcommand{\matchinglabelsdotted}[1]{
    \foreach \i/\j in {#1} {
        \arclabeldotted{\i}{\j}{\i}{\j};
    };
}

\newcommand{\matchinglabelsgray}[1]{
    \foreach \i/\j in {#1} {
        \arclabelgray{\i}{\j}{\i}{\j};
    };
}

\newcommand{\matchingperm}[1]{
    \pgfmathsetmacro\n{dim(#1)};
    % Now \n stores the number of entries of the permutation.
    \foreach \j [count=\i] in {#1} {
        \arclabel{\i}{{2*\n+1-\j}}{\i}{\j};
    };
}

\newcommand{\matchingshadeedge}[1]{
    % Just like drawing an edge, e.g., 3/6
    \foreach \i/\j in {#1} {
        \draw[color=lightgray, fill=lightgray] ({\i-0.25},0)
            arc (180:0:{(\j-\i)/2 + 0.25})
            arc (0:-180:0.25)
            arc (0:180:{(\j-\i)/2 - 0.25})
            arc (0:-180:0.25)--cycle;   
    };
}

\newcommand{\matchingshadeendpoints}[1]{
    % Multiple groups can be input, like 1/4, 5/8
    \foreach \i/\j in {#1} {
        \draw[color=lightgray, fill=lightgray] (\i,-0.25)
            arc (270:90:0.25)
            --(\j,0.25)
            arc (90:-90:0.25)
            --(\i,-0.25);
    };
}

\newcommand{\matchingshadegroup}[4]{
    % Give two groups, like {1}{4}{5}{8}. Will shade whole region of edges from vertices 1-4 to vertices 5-8.
    \draw[color=lightgray, fill=lightgray] ({#1-0.25},0)
        arc (180:0:{(#4-#1)/2 + 0.25}) % outer arc
        arc (0:-90:0.25)
        --(#3,-0.25) % right-hand bottom
        arc (270:180:0.25)
        arc (0:180:{(#3-#2)/2 - 0.25}) % inner arc
        arc (0:-90:0.25)
        --(#1,-0.25) % left-hand bottom
        arc (270:180:0.25)--cycle;
}

\newcommand{\matchinggroupdark}[4]{
    % Give two groups, like {1}{4}{5}{8}. Will shade whole region of edges from vertices 1-4 to vertices 5-8.
    \draw ({#1-0.25},0)
        arc (180:0:{(#4-#1)/2 + 0.25}) % outer arc
        arc (0:-90:0.25)
        --(#3,-0.25) % right-hand bottom
        arc (270:180:0.25)
        arc (0:180:{(#3-#2)/2 - 0.25}) % inner arc
        arc (0:-90:0.25)
        --(#1,-0.25) % left-hand bottom
        arc (270:180:0.25)--cycle;
}

\newcommand{\matchingshadegroupdark}[4]{
    % Give two groups, like {1}{4}{5}{8}. Will shade whole region of edges from vertices 1-4 to vertices 5-8.
    \draw[color=black, fill=lightgray] ({#1-0.25},0)
        arc (180:0:{(#4-#1)/2 + 0.25}) % outer arc
        arc (0:-90:0.25)
        --(#3,-0.25) % right-hand bottom
        arc (270:180:0.25)
        arc (0:180:{(#3-#2)/2 - 0.25}) % inner arc
        arc (0:-90:0.25)
        --(#1,-0.25) % left-hand bottom
        arc (270:180:0.25)--cycle;
}

%
%
%
%

% Drawing (small) matchings (all have width = 0.1*(verts-1) and height=0.15

\newcommand{\matchnest}{
    \begin{tikzpicture}[scale=.1, yscale=1.5, anchor=base]
        \arcskinnyplain{1}{4};
        \arcskinnyplain{2}{3};
    \end{tikzpicture}
}

\newcommand{\matchcross}{
    \begin{tikzpicture}[scale=.1, yscale=2.25, anchor=base]
        \arcskinnyplain{1}{3};
        \arcskinnyplain{2}{4};
    \end{tikzpicture}
}

\newcommand{\matchind}{
    \begin{tikzpicture}[scale=.1, yscale=4.5, anchor=base]
        \arcskinnyplain{1}{2};
        \arcskinnyplain{3}{4};
    \end{tikzpicture}
}

\newcommand{\matchsmall}[1]{
    \begin{tikzpicture}[scale=.1, anchor=base]
        % We need to read through so that we can rescale y so that the height is precisely 0.15
        \def\h{0};
        \def\maxh{0};
        \foreach \i/\j in {#1} {
            \pgfmathparse{\j-\i};
            \let\h\pgfmathresult;
            \pgfmathifthenelse{\h>\maxh}{\h}{\maxh};
            \global\let\maxh\pgfmathresult;
        };
        \pgftransformyscale{{4.5/\maxh}};
        \foreach \i/\j in {#1} {
            \arcskinnyplain{\i}{\j};
        };
    \end{tikzpicture}
}

\newcommand{\matchpermsmall}[1]{
    \begin{tikzpicture}[scale=.1, anchor=base]
        \pgfmathsetmacro\n{dim(#1)};
        % Now \n stores the length of the permutation.
        % We need to read through again so that we can rescale y so that the height is precisely 0.15
        \def\h{0};
        \def\maxh{0};
        \foreach \j [count=\i] in {#1} {
            \pgfmathparse{2*\n+1-\j-\i};
            \let\h\pgfmathresult;
            \pgfmathifthenelse{\h>\maxh}{\h}{\maxh};
            \global\let\maxh\pgfmathresult;
        };
        \pgftransformyscale{{4.5/\maxh}};
        \foreach \j [count=\i] in {#1} {
            \arcskinnyplain{\i}{{2*\n+1-\j}};
        };
    \end{tikzpicture}
}

%
%
%
%

% Plotting partitions.
\newcommand{\plotpartition}[1]{
    \foreach \i [count=\j] in {#1} {
        \draw [thick, line cap=round] (0,-\j)--(\i,-\j)--(\i,-\j-1)--(0,-\j-1)--cycle;
        \foreach \k in {1,2,...,\i}{
            \draw [thick, line cap=round] ({\k-1},-\j)--({\k-1},{-\j-1});
        };
    };
}
\newcommand{\plotpartitiongray}[1]{
    \foreach \i [count=\j] in {#1} {
        \draw [lightgray, thick, line cap=round] (0,-\j)--(\i,-\j)--(\i,-\j-1)--(0,-\j-1)--cycle;
        \foreach \k in {1,2,...,\i}{
            \draw [lightgray, thick, line cap=round] ({\k-1},-\j)--({\k-1},{-\j-1});
        };
    };
}

% Example:
%       \begin{center}
%           \begin{tikzpicture}[scale=0.25]
%               \plotpartitiongray{5,3,2,1,6,1,1};
%               \plotpartition{4,1,1,1,1};
%           \end{tikzpicture}
%       \end{center}

%
%
%
%

% Plotting compositions (skyline diagrams).
\newcommand{\plotskyline}[1]{
    \foreach \j [count=\i] in {#1} {
        \draw [thick, line cap=round] (\i-1,0)--(\i-1,\j)--(\i,\j)--(\i,0)--cycle;
        \foreach \k in {1,2,...,\j}{
            \draw [thick, line cap=round] (\i-1,\k-1)--(\i,\k-1);
        };
    };
}
\newcommand{\plotskylinegray}[1]{
    \foreach \j [count=\i] in {#1} {
        \draw [lightgray, thick, line cap=round] (\i-1,0)--(\i-1,\j)--(\i,\j)--(\i,0)--cycle;
        \foreach \k in {1,2,...,\j}{
            \draw [lightgray, thick, line cap=round] (\i-1,\k-1)--(\i,\k-1);
        };
    };
}
\newcommand{\plotskylineshaded}[1]{
    \foreach \j [count=\i] in {#1} {
        \draw [lightgray, fill=lightgray, thick, line cap=round] (\i-1,0)--(\i-1,\j)--(\i,\j)--(\i,0)--cycle;
        \foreach \k in {1,2,...,\j}{
            \draw [lightgray, fill=lightgray, thick, line cap=round] (\i-1,\k-1)--(\i,\k-1);
        };
    };
}

% Example:
%       \begin{center}
%           \begin{tikzpicture}[scale=0.25]
%               \plotskylinegray{5,3,2,4,6,1,2};
%               \plotskyline{4,1,1,1,5,1,2};
%           \end{tikzpicture}
%       \end{center}

%
%
%
%

%%%%%%%%%%%%%%%%%%%%%%%%%%%%%%%%%%%%%%%%%%%%%%%%%%%%%%%%%%%%%%%%
%
% Code for inline 1x2 (vert) and 2x1 (horiz) grids. June 2017.
%
%%%%%%%%%%%%%%%%%%%%%%%%%%%%%%%%%%%%%%%%%%%%%%%%%%%%%%%%%%%%%%%%

\newcommand{\onegrid}[1]{
\begin{tikzpicture}[scale=1, baseline = 1pt]
    \pgftransformxscale{0.25};
    \pgftransformyscale{0.25};
    \draw [semithick, line cap=round] (0,0) rectangle (1,1);
    \ifthenelse{#1>0}{
        \draw [semithick, line cap=round] (0,0)--(1,1);
    }{
        \draw [semithick, line cap=round] (0,1)--(1,0);
    };
\end{tikzpicture}
}

\newcommand{\onegriddotNW}{
\begin{tikzpicture}[scale=1, baseline = 1pt]
    \pgftransformxscale{0.25};
    \pgftransformyscale{0.25};
    \draw [draw=black, fill=black] (0,1) circle (4pt);
    \draw [semithick, line cap=round] (0,0) rectangle (1,1);    
    \draw [semithick, line cap=round] (0,0)--(1,1);
\end{tikzpicture}
}

\newcommand{\onegriddotSE}{
\begin{tikzpicture}[scale=1, baseline = 1pt]
    \pgftransformxscale{0.25};
    \pgftransformyscale{0.25};
    \draw [draw=black, fill=black] (1,0) circle (4pt);
    \draw [semithick, line cap=round] (0,0) rectangle (1,1);    
    \draw [semithick, line cap=round] (0,0)--(1,1);
\end{tikzpicture}
}

\newcommand{\gridsmallhoriz}[1]{
% Parameters should be +1 or -1, such as \gridsmallhoriz{1,-1,-1}.
    \begin{tikzpicture}[scale=1, baseline = 1pt]
        \def\gridheight{1};
        \def\gridwidth{1}
        % \foreach \dir [count=\gridwidth] in {#1} {
        %     % We read through first to determine how wide the grid is.
        % };
        \pgftransformxscale{{0.25/\gridwidth}};
        \pgftransformyscale{{0.25/\gridheight}};
        \foreach \dir [count=\i] in {#1} {
            % Now we have to draw the lines. We start at (0,0) or (0,1).
            \ifthenelse{\dir>0}{
                \draw [semithick, line cap=round] ({\i-1}, 0)--(\i, 1);
            }{
                \draw [semithick, line cap=round] ({\i-1}, 1)--(\i, 0);
            };
            \draw [semithick, line cap = round] ({\i-1}, 0) rectangle ({\i},1);
        };
    \end{tikzpicture}
}

\newcommand{\gridhoriz}[1]{
 % Parameters should be +1 or -1, such as \gridsmallhoriz{1,-1,-1}.
\begin{tikzpicture}[scale=1, anchor=base]
 \pgftransformxscale{0.225};
\pgftransformyscale{0.225};
   \foreach \dir [count=\i] in {#1} {
     % Now we have to draw the lines. We start at (0,0) or (0,1).
        \ifthenelse{\dir>0}{
            \draw [semithick, line cap=round] ({\i-1}, 0)--(\i, 1);
        }{
            \draw [semithick, line cap=round] ({\i-1}, 1)--(\i, 0);
        };
    \draw [semithick, line cap = round] ({\i-1}, 0) rectangle ({\i},1);
   };
 \end{tikzpicture}
}

\newcommand{\gridhorizframe}[1]{
 % Parameters should be +1 or -1, such as \gridsmallhoriz{1,-1,-1}.
\begin{tikzpicture}[scale=1, anchor=base]
 \pgftransformxscale{0.225};
\pgftransformyscale{0.225};
   \foreach \dir [count=\i] in {#1} {
     % Now we have to draw the lines. We start at (0,0) or (0,1).
     \ifthenelse{\dir>0}{
       \draw [semithick, line cap=round] ({\i-1}, 0)--(\i, 1);
     }{
       \draw [semithick, line cap=round] ({\i-1}, 1)--(\i, 0);
     };
     % Draw the frame.
     \draw [semithick, line cap = round] ({\i-1}, 0) rectangle ({\i},1);
   };
 \end{tikzpicture}
}

\newcommand{\gridsmallvert}[1]{
 % Parameters should be +1 or -1, such as \gridsmallhoriz{1,-1,-1}.
 % These are read from bottom to top.
\begin{tikzpicture}[scale=1, anchor=base]
   \def\gridwidth{1};
   \foreach \dir [count=\gridheight] in {#1} {
     % We read through first to determine how wide the grid is.
   };
 \pgftransformxscale{{0.225/\gridwidth}};
\pgftransformyscale{{0.225/\gridheight}};
   \foreach \dir [count=\i] in {#1} {
     % Now we have to draw the lines. We start at (0,0) or (0,1).
     \ifthenelse{\dir>0}{
   \draw [semithick, line cap=round] (0, {\i-1})--(1, \i);
     }{
       \draw [semithick, line cap=round] (0, \i)--(1, {\i-1});
     };
   };
 \end{tikzpicture}
}

\newcommand{\gridverysmallvert}[1]{
 % Parameters should be +1 or -1, such as \gridsmallhoriz{1,-1,-1}.
 % These are read from bottom to top.
\begin{tikzpicture}[scale=1, anchor=base]
   \def\gridwidth{1};
   \foreach \dir [count=\gridheight] in {#1} {
     % We read through first to determine how wide the grid is.
   };
 \pgftransformxscale{{0.175/\gridwidth}};
 \pgftransformyscale{{0.175/\gridheight}};
   \foreach \dir [count=\i] in {#1} {
     % Now we have to draw the lines. We start at (0,0) or (0,1).
     \ifthenelse{\dir>0}{
   \draw [semithick, line cap=round] (0, {\i-1})--(1, \i);
     }{
       \draw [semithick, line cap=round] (0, \i)--(1, {\i-1});
     };
   };
 \end{tikzpicture}
}

\newcommand{\gridsmallhorizfn}[1]{
    % For use in footnotes (where everything has to be smaller, about 80%)
    % Parameters should be +1 or -1, such as \gridsmallhoriz{1,-1,-1}.
    \begin{tikzpicture}[scale=1, anchor=base]
    \def\gridheight{1};
    \foreach \dir [count=\gridwidth] in {#1} {
        % We read through first to determine how wide the grid is.
    };
    \pgftransformxscale{{0.155/\gridwidth}};
    \pgftransformyscale{{0.155/\gridheight}};
    \foreach \dir [count=\i] in {#1} {
        % Now we have to draw the lines. We start at (0,0) or (0,1).
        \ifthenelse{\dir>0}{
            \draw [semithick, line cap=round] ({\i-1}, 0)--(\i, 1);
        }{
            \draw [semithick, line cap=round] ({\i-1}, 1)--(\i, 0);
        };
    };
  \end{tikzpicture}
}

\newcommand{\gridsmalltwobyvert}[2]{
    % Parameters should be +1 or -1, such as \gridsmalltwobyvert{1,-1,-1}.
    % These are read from bottom to top.
    % In this version, we have TWO COLUMNS.
    \begin{tikzpicture}[scale=1, anchor=base]
    \def\gridwidth{2};
    \foreach \dir [count=\gridheight] in {#1} {
        % We read through first to determine how wide the grid is.
    };
    \pgftransformxscale{{0.225/\gridwidth}};
    \pgftransformyscale{{0.225/\gridheight}};
    \foreach \dir [count=\i] in {#1} {
        % First column.
        \ifthenelse{\dir=1 \OR \dir=-1}{
            \ifthenelse{\dir>0}{
                \draw [semithick, line cap=round] (0, {\i-1})--(1, \i);
            }{
                \draw [semithick, line cap=round] (0, \i)--(1, {\i-1});
            };
        }{};
    };
    \foreach \dir [count=\i] in {#2} {
        % Second column.
        \ifthenelse{\dir=1 \OR \dir=-1}{
            \ifthenelse{\dir>0}{
                \draw [semithick, line cap=round] (1, {\i-1})--(2, {\i  });
            }{
                \draw [semithick, line cap=round] (1, {\i  })--(2, {\i-1});
            };
        }{};
    };
    \end{tikzpicture}
}

\newcommand{\gridsmalltwobyvertframe}[2]{
    % Parameters should be +1 or -1, such as \gridsmalltwobyvert{1,-1}{-1,0}.
    % These are read from bottom to top.
    % In this version, we have TWO COLUMNS.
    \begin{tikzpicture}[scale=1, anchor=base]
    \def\gridwidth{2};
    \def\gridheight{2};
    % \foreach \dir [count=\gridheight] in {#1} {
    %     % We read through first to determine how wide the grid is.
    % };
    \pgftransformxscale{{0.25/\gridwidth}};
    \pgftransformyscale{{0.25/\gridheight}};
    \foreach \dir [count=\i] in {#1} {
        % First column.
        \ifthenelse{\dir=1 \OR \dir=-1}{
            \ifthenelse{\dir>0}{
                \draw [semithick, line cap=round] (0, {\i-1})--(1, \i);
            }{
                \draw [semithick, line cap=round] (0, \i)--(1, {\i-1});
            };
        }{};
        \draw [line cap = round] (0, {\i-1}) rectangle (1, {\i});
    };
    \foreach \dir [count=\i] in {#2} {
        % Second column.
        \ifthenelse{\dir=1 \OR \dir=-1}{
            \ifthenelse{\dir>0}{
                \draw [semithick, line cap=round] (1, {\i-1})--(2, {\i  });
            }{
                \draw [semithick, line cap=round] (1, {\i  })--(2, {\i-1});
            };
        }{};
        \draw [line cap = round] (1, {\i-1}) rectangle (2, {\i});
    };
    \end{tikzpicture}
}

\newcommand{\gridsmallthreebyvertframe}[3]{
    % Parameters should be +1 or -1, such as \gridsmallthreebyvert{1,0,-1}{0,-1,1}{-1,1,0}.
    % These are read from bottom to top.
    % In this version, we have THREE COLUMNS.
    \begin{tikzpicture}[scale=1, baseline=1pt]
    \def\gridwidth{2};
    \def\gridheight{2};
    % \foreach \dir [count=\gridheight] in {#1} {
    %     % We read through first to determine how wide the grid is.
    % };
    \pgftransformxscale{{0.25/\gridwidth}};
    \pgftransformyscale{{0.25/\gridheight}};
    \foreach \dir [count=\i] in {#1} {
        % First column.
        \ifthenelse{\dir=1 \OR \dir=-1}{
            \ifthenelse{\dir>0}{
                \draw [semithick, line cap=round] (0, {\i-1})--(1, \i);
            }{
                \draw [semithick, line cap=round] (0, \i)--(1, {\i-1});
            };
        }{};
        \draw [line cap = round] (0, {\i-1}) rectangle (1, {\i});
    };
    \foreach \dir [count=\i] in {#2} {
        % Second column.
        \ifthenelse{\dir=1 \OR \dir=-1}{
            \ifthenelse{\dir>0}{
                \draw [semithick, line cap=round] (1, {\i-1})--(2, {\i  });
            }{
                \draw [semithick, line cap=round] (1, {\i  })--(2, {\i-1});
            };
        }{};
        \draw [line cap = round] (1, {\i-1}) rectangle (2, {\i});
    };
    \foreach \dir [count=\i] in {#3} {
        % Second column.
        \ifthenelse{\dir=1 \OR \dir=-1}{
            \ifthenelse{\dir>0}{
                \draw [semithick, line cap=round] (2, {\i-1})--(3, {\i  });
            }{
                \draw [semithick, line cap=round] (2, {\i  })--(3, {\i-1});
            };
        }{};
        \draw [line cap = round] (2, {\i-1}) rectangle (3, {\i});
    };
    \end{tikzpicture}
}

